\documentclass[preprint,12pt]{elsarticle}
\usepackage{amsmath,amssymb,graphicx,booktabs,hyperref,url}
\usepackage{multirow}
\usepackage{array}

\journal{Journal of Machine Learning Research}

\begin{document}

\begin{frontmatter}

\title{Fragile Diagnostics: Leave-One-Out Ensemble Selection is Seed-Sensitive, Detector-Dependent, and Universally Unreliable}

\author[aff1]{Md. Omar Faruk Rabby}
\affiliation[aff1]{organization={Department of Computer Science and Engineering, Southeast University},
                     city={Dhaka},
                     postcode={1208},
                     country={Bangladesh}}

\begin{abstract}
% Written in STEP P5
\end{abstract}

\begin{keyword}
anomaly detection \sep ensemble pruning \sep leave-one-out contribution \sep network intrusion detection \sep reproducibility
\end{keyword}

\end{frontmatter}

%=============================================================================
\section{Introduction}
%=============================================================================

% Paragraph 1: NIDS ensemble construction practical problem.
Network intrusion detection systems (NIDS) increasingly employ ensemble methods to combine multiple anomaly detectors, each capturing different aspects of attack behavior. The ensemble construction problem---selecting which detectors to include---is critical: including harmful detectors degrades performance, while excluding useful ones reduces coverage. Leave-one-out (LOO) contribution, which measures each detector's marginal impact by comparing ensemble performance with and without it, is the dominant selection heuristic in this domain.

% Paragraph 2: Gap---no ground-truth validation.
Despite widespread adoption, LOO's diagnostic reliability has never been validated against known ground truth. The fundamental assumption---that a negative LOO contribution indicates a harmful detector---conflates two distinct failure modes: genuine detector harm (the detector actively degrades ensemble performance) and detector redundancy (the detector provides no new information, causing its removal to have negligible or negative effect on a different metric). This conflation makes LOO unreliable for automated ensemble pruning.

% Paragraph 3: Three contributions.
We present three contributions. First, a controlled synthetic benchmark with 2,000 ensemble configurations where detector roles (harmful, redundant, beneficial) are known, revealing that LOO detects harmful detectors with perfect recall (1.000) but misclassifies 67\% of redundant detectors as beneficial (recall 0.327). Second, a multi-dataset evaluation across 3 NIDS datasets, 3 categorical encoders, and 20 seeds (180 configurations) demonstrating that sign-conflict between LOO's validation signal and held-out test effect occurs at 48\%, independent of ensemble correlation structure. Third, epsilon-VRG, a practical guardrail that reduces redundancy false-pruning by 66\% without retaining harmful detectors.

% Paragraph 4: Practical implication.
Our findings do not advocate abandoning LOO as a diagnostic signal. Rather, we show that LOO's directional claims are inherently seed-sensitive and detector-dependent, requiring complementary guardrails for reliable ensemble construction. We provide practical guidelines for practitioners building NIDS ensembles.

% Paragraph 5: Roadmap.
The remainder of this paper is organized as follows. Section~\ref{sec:related} reviews related work. Section~\ref{sec:formulation} presents the problem formulation. Section~\ref{sec:phase2} describes the controlled ground-truth study. Section~\ref{sec:phase3} presents the real-data multi-dataset evaluation. Section~\ref{sec:fragility} analyzes seed fragility. Section~\ref{sec:discussion} discusses implications and limitations. Section~\ref{sec:guidelines} provides practical guidelines. Section~\ref{sec:conclusion} concludes.

%=============================================================================
\section{Related Work}
\label{sec:related}
%=============================================================================

\subsection{Ensemble pruning for anomaly detection}
% 5-8 sentences covering ensemble pruning literature, focusing on NIDS applications.
% Key references: boosting, bagging, selective ensemble, ECOC.
Placeholder: Review ensemble pruning methods in anomaly detection, emphasizing the dominance of LOO-based selection in NIDS literature.

\subsection{Leave-one-out and marginal contribution}
% 5-8 sentences on LOO theory, Jackknife, influence functions.
% Key distinction: LOO in supervised vs unsupervised settings.
Placeholder: Discuss LOO contribution computation, Jackknife methods, and the difference between supervised and unsupervised LOO interpretation.

\subsection{Leakage in unsupervised ML}
% 3-5 sentences on how categorical encoding can leak labels.
Placeholder: Address the under-appreciated risk of label leakage through categorical encoding in unsupervised anomaly detection pipelines.

\subsection{NIDS ensemble construction}
% 3-5 sentences on practical NIDS ensemble construction challenges.
Placeholder: Review how existing NIDS frameworks construct ensembles and where LOO selection fits in practice.

%=============================================================================
\section{Problem Formulation}
\label{sec:formulation}
%=============================================================================

\subsection{Rank-average ensemble}
% Define the ensemble: score = rank average of K detectors.
% Q_all = ensemble AUROC on validation set.

\subsection{Leave-one-out contribution}
% C_i = Q_all - Q_{-i}
% Interpretation: positive = detector helps, negative = detector hurts.

\subsection{TestEffect and sign conflict}
% TE_i = test_AUROC_all - test_AUROC_{-i}
% Sign conflict: sign(C_i) != sign(TE_i)
% Material conflict: sign(C_i) != sign(TE_i) AND |TE_i| > threshold.

\subsection{Epsilon-VRG}
% epsilon = alpha * AUROC_full(D_val)
% Retain if C_i > -epsilon, exclude if C_i <= -epsilon.
% alpha = 0.01 (pre-specified).

%=============================================================================
\section{Controlled Ground-Truth Study (Phase 2)}
\label{sec:phase2}
%=============================================================================

\subsection{Synthetic benchmark design}
% 200 configurations x 10 seeds = 2000 runs.
% Detector roles: harmful (H), redundant (R), beneficial (B).
% Role assignment based on known AUROC and correlation structure.

\subsection{Hypotheses}

\begin{itemize}
\item H1: LOO detects harmful detectors (recall $\geq 0.80$).
\item H2: LOO prunes more redundant than beneficial detectors (R recall $>$ B recall).
\item H3: Sign-conflict rate correlates with ensemble correlation (Spearman $\rho > 0.3$).
\item H4: epsilon-VRG reduces redundancy false-pruning by $\geq 50\%$.
\end{itemize}

\subsection{Results}
% H1: PASS (recall 1.000, n=225)
% H2: PASS (R recall 0.327, B recall 0.819)
% H3: REJECTED (rho ~ 0)
% H4: PASS (R-fp -66%, H retention 0, B sl -62%)

\subsection{Discussion}
% LOO is a good harmful-detector detector but poor redundant-detector discriminator.
% Sign-conflict is noise-driven, not correlation-dependent.

%=============================================================================
\section{Real-Data Multi-Dataset Evaluation (Phase 3)}
\label{sec:phase3}
%=============================================================================

\subsection{Experimental design}
% 3 datasets x 3 encoders x 20 seeds = 180 configs.
% 6 detectors (IF, LOF, OCSVM, AE, ECOD, HBOS).
% Methods: EW, NL, EV.

\subsection{Results}
% Per-dataset method means table.
% Holm K=9 corrected comparisons.

\subsection{Statistical analysis}
% Wilcoxon signed-rank, Holm correction, hierarchical model for seed variance.

%=============================================================================
\section{Seed-Fragility Analysis (Phase 3 Meta-Finding)}
\label{sec:fragility}
%=============================================================================

\subsection{Direction flips across seeds}
% EW > NL in only 4/20 seeds (NSL-KDD Phase 1 replication).
% 16/20 seeds show NL > EW.

\subsection{Root causes}
% (1) Detector set effect: ECOD/HBOS in 6-detector set degrades EW.
% (2) Seed fragility: direction flips across seeds.
% Pipeline parity confirmed: 4-detector subset shows EW diff < 0.001.

\subsection{Implications}
% Phase 1's encoder-dependent EW > NL reversal does NOT replicate.
% H5 reframed as seed-fragility finding.

%=============================================================================
\section{Discussion}
\label{sec:discussion}
%=============================================================================

\subsection{Why LOO conflates redundancy with harm}
% Mathematical explanation: removal of a redundant detector changes rank normalization.
% LOO cannot distinguish "removes useful information" from "changes rank distribution."

\subsection{Practical impact}
% For NIDS practitioners: LOO alone is insufficient for reliable selection.
% Need complementary guardrails (epsilon-VRG, sign-conflict monitoring).

\subsection{Limitations}
% 1. Only 3 NIDS datasets.
% 2. Only 6 detector families.
% 3. K=9 is post-hoc.
% 4. AE failure may be architecture-specific.
% 5. No cross-dataset transfer evaluation.

%=============================================================================
\section{Practical Guidelines for Practitioners}
\label{sec:guidelines}
%=============================================================================

% Table with scenarios and recommended actions.
% Content written in STEP P6.

%=============================================================================
\section{Conclusion}
\label{sec:conclusion}
%=============================================================================

We have demonstrated that leave-one-out ensemble selection, despite widespread adoption, is fundamentally unreliable as a standalone diagnostic tool. Our controlled ground-truth study reveals that LOO conflates redundancy with harm (R recall 0.327), and our multi-dataset evaluation shows that its directional claims are seed-sensitive (sign-conflict rate 48\%) and detector-dependent. We propose epsilon-VRG as a practical guardrail and provide guidelines for practitioners. All code and data are publicly available for reproducibility.

%=============================================================================
% References
%=============================================================================
\bibliographystyle{elsarticle-num}
\bibliography{references}

\end{document}
